\section*{\S~5b Dienstgradschlüssel}
\label{sec:dienstgradschluessel}
\addcontentsline{toc}{section}{\S~5b Dienstgradschlüssel}

\begin{enumerate}[label=(\arabic*)]
	\item Da je nach Mitgliederanzahl bestimmte Ränge nicht benötigt werden, gilt der im Folgenden aufgelistete Dienstgradschlüssel. Der Dienstgradschlüssel bezieht sich immer nur auf Dienstgrade, die von aktiven Mitgliedern besetzt werden.
        \item Wenn eine unten beschriebene Schwelle für einen Dienstgrad erreicht wurde, so entscheidet die Leitung, wann die für eine Beförderung erforderlichen Maßnahmen ergriffen werden. Wird die Schwelle für eine solche Beförderung vor Durchführung dieser Beförderung wieder unterschritten, wird die Beförderung nicht durchgeführt. Wird die Schwelle für eine solche Beförderung nach Durchführung dieser Beförderung wieder unterschritten, bleibt die Beförderung bestehen. Bei erneutem Erreichen dieser Schwelle wird nicht erneut befördert.
\end{enumerate}

\begin{center}
\begin{tabular}{|p{8cm}|p{6cm}|}
	\hline
	\textbf{Mitgliederanzahl (ohne Leitung)} & \textbf{Dienstgrade} \\
	\hline
	Weniger als 8 aktive Mitglieder & Mannschaftsdienstgrade bis Oberstabsgefreiter + Dienstgrade der Leitung \\
	\hline
	Mindestens 8 aktive Mitglieder & Vorhergehende + einen Korporal/Stabskorporal \\
	\hline
	Mindestens 9 aktive Mitglieder & Vorhergehende + einen Unteroffizier ohne Portepee \\
	\hline
	Mindestens 13 aktive Mitglieder & Vorhergehende + einen Korporal/Stabskorporal und einen Unteroffizier mit Portepee \\
	\hline
	Mindestens 18 aktive Mitglieder & Vorhergehende + einen Korporal/Stabskorporal und einen Unteroffizier ohne Portepee \\
	\hline
	Mindestens 30 aktive Mitglieder & Vorhergehende + zwei Unteroffiziere ohne Portepee und zwei Unteroffiziere mit Portepee \\
	\hline
\end{tabular}
\end{center}
